\subsection{在作用下稳定的子集。诱导作用}

定义2. 一个子集 \(\mathbf{A}\) （属于一个集合 \(\mathbf{E}\) ）称为在某个作用 \(\alpha \mapsto f_{\alpha}\) （ \(\Omega\) 在 \(\mathbf{E}\) 上的作用）下稳定，若 \(f_{\alpha}(\mathbf{A}) \subset \mathbf{A}\) （对所有 \(\alpha \in \Omega\) ）。一个元素 \(x\) （属于 \(\mathbf{E}\) ）称为在元素 \(\alpha\) （属于 \(\Omega\) ）下不变，若 \(f_{\alpha}(x) = x\) .

在给定作用下，E 的稳定子集族的交集是稳定的。因此，存在包含 E 的给定子集 X 的最小稳定子集；称它为由 X 生成；它由形如 \((f_{\alpha_{1}} \circ f_{\alpha_{2}} \circ \cdots \circ f_{\alpha_{n}})(x)\) 的元素（其中 \(x \in X, n \geqslant 0, \alpha_{1} \in \Omega\) 对所有 i 成立）组成。

注。设 E 是一个律记为 \(\top\) 的广群。应当指出，E 的子集 A 在 E 对自身的左作用下稳定，并不一定在 E 对自身的右作用下稳定；E 的在自身左（相应地，右）作用下稳定的子集 A 在 E 的律下稳定，但反之一般不成立。更确切地说，A 在 E 的律下稳定当且仅当 \(A \top A \subset A\) ，而 A 在 E 对自身的左（相应地，右）作用下稳定当且仅当 \(E \top A \subset A\) （相应地 \(A \top E \subset A\) ).

示例。取广群 E 为带乘法的集合 N。集合 \(\{1\}\) 在 N 的内部律下是稳定的，但在 N 对自身的作用下由 \(\{1\}\) 生成的稳定子集是整个 N。

定义3. 设 \(\alpha \mapsto f_{\alpha}\) 是 \(\Omega\) 在 E 上的一个作用，且 A 是 E 的一个稳定子集。将元素 \(\alpha \in \Omega\) 与 \(f_{\alpha}\) 在 A 上的限制（视为 A 到自身的映射）相关联的映射是 \(\Omega\) 在 A 上的一个作用，称为由给定作用诱导的作用。

